\section{Limits of the estimates and a projective comparison}
\subsection{The dimension ceiling in four weights}
Let \(J_{p,t}\) evaluate \(P(X,y_1,\ldots,y_4)\) at the four finite
polylogarithms modulo \(X^p\), with \(\deg_X P\le t^2\) and
\(\deg_y P\le t\) in total. Its source has dimension
\(D=(t^2+1)\binom{t+4}{4}\sim t^6/24\). An injectivity theorem up
to \(t=\lfloor cp^{1/6}\rfloor\) for a fixed small \(c>0\) would give
\(O(p^{5/6})\) in Lemma~\ref{lem:mixed-disc-count} for \(n=4\), by
leaving the auxiliary polynomial and multiplicity steps unchanged.
The example in Theorem~\ref{thm:zero} explains why an improvement of this
kind must exploit the particular system of finite polylogarithms.

For \(t=1,2,3\) one can prove maximal rank for every prime by an integer
calculation. On coefficients of index at most \(D\), the Taylor coefficients
of the characteristic-zero series \(\sum_{j\ge1}X^j/j^i\) specialise
to the finite series for every \(p>D\). The pure \(X\)-power columns give
an identity block on rows \(0,\ldots,t^2\). Complete the block in two ways
by using either the next \(D-t^2-1\) rows or the same rows shifted up
by one index. Clear denominators in each remaining row. The clearing
factors are units for \(p>D\). If both integer minors vanished modulo such a
prime it would divide their gcd. The following table gives the bit lengths
of the minors, their gcd after removing all factors at most \(D\), and
the number of smaller primes checked directly for surjectivity.
\begin{center}
\begin{tabular}{rrrrr}
\toprule
\(t\)&\(D\)&bit lengths&remaining gcd&primes \(p\le D\)\\
\midrule
1&10&37,\ 37&1&4\\
2&75&9636,\ 9869&1&21\\
3&350&332544,\ 334400&1&70\\
\bottomrule
\end{tabular}
\end{center}
Thus \(\operatorname{rank}J_{p,t}=\min(D,p)\) for these three values of \(t\).
The second minor matters already at \(t=2,p=109\), where the first 75
rows have rank 74. In boundary tests up to \(t=7\), the calculation gives
ranks \(16493\) and \(16500\) at \(p=16493\) and \(16519\) respectively
when \(D=16500\). A bound for growing \(t\) remains open.

\subsection{Density one at depth four}
For \(S=\{2,3\}\), equality with the 21 integral points for a density-one
set of auxiliary primes would require excluding extra common zeros on
both refinement components. Even on \((1,1)\), where the logarithm
restricts the points to roots of unity, two difficulties remain. The
possible orders grow with \(p-1\). At a fixed order three or four the
conditions at weights two and four evaluate different branches
\(\chi\omega^{-1}\) and \(\chi\omega^{-3}\) of the \(p\)-adic \(L\)-function.
The calculation at either conductor in Section~\ref{sec:exceptional}
rules out individual primes by finding a nonzero digit. It gives no
density estimate for simultaneous vanishing as \(p\) varies.
Betts, Kumpitsch and L\"udtke isolate \(-1\) using a weight depending on
\(p\) \cite[Proposition~5.15]{BKL}. The question here fixes the depth.

\subsection{A genus-three calculation and the Cartier operator}
Consider the smooth projective model \(C\) of
\[
 y^2=f(x)=4x^7+9x^6-8x^5-36x^4-16x^3+32x^2+32x+8.
\]
Its Jacobian \(J\) has rank one. The divisor \(D=(-1,-1)-\infty\) has
infinite order \cite[Example~4.1]{GenusThree}. At a good odd prime we
count the discs meeting the inverse image of \(\mathbb Q_p\log D\)
under \(P\mapsto\log[P-\infty]\). This definition uses the rational
Mordell--Weil span and requires no identification with a Selmer group.

Away from branch points choose a lift \(Q\) with local coordinate
\(x=x(Q)+pw\). Project the integrals of \(x^i\,dx/(2y)\), \(0\le i\le2\),
onto an integral basis of the two-dimensional annihilator of \(\log D\).
Write \(c_1,c_2\) for the two projected constants and \(u_1,u_2\) for their
tangent coefficients modulo \(p\). A lift requires \(c_i\in p\mathbb Z_p\)
and
\[
                 c_i/p+u_iw\equiv 0 \pmod p,\qquad i=1,2.
\]
For a nonzero slope vector, compatibility amounts to
\((c_1/p)u_2-(c_2/p)u_1\equiv0\pmod p\). When both slopes vanish, both
reduced constants must vanish. A change of lift translates the constant
vector along the slope vector. Stoll also uses the direction of a
normalised abelian logarithm, combining it at \(p=2\) with a global
Selmer-group test for rational points \cite[Proposition~2.8]{Stoll}.

Table~\ref{tab:genus} gives the calculation for all 45 odd primes below
200. For \(n=\#J(\mathbb F_p)\), compute \([n](Q-\infty)\) in the
formal neighbourhood of the identity. Translating by a regular
degree-three divisor allows Cantor arithmetic on a chart with simple
support. Integrating between the two sets of nearby support points and
taking the trace gives \(p^{-1}\log([n](Q-\infty))\bmod p^3\).
For \(p\ge5\) the omitted terms of degree \(j\ge4\) have valuation at least
\(j-v_p(j)\ge4\) before dividing by \(p\). Testing all admissible local
digits gives the upper count. Lower witnesses are the Weierstrass
points and the available pairs above \(x=-1,1,0\). The former two
pairs are rational. For \(x=0\) a torsion witness follows from
\[
 \begin{aligned}
 B &=x^2+4x+2,\\
 A &=-x^6-16x^5-36x^4+72x^2+64x+16,&
 A^2-8fB^2&=x^{12}.
 \end{aligned}
\]
At \(P=(0,2\sqrt2)\) the function \(A-2\sqrt2\,yB\) has a zero
of order 12 and a unique pole at infinity. Its divisor proves that
\(12[P-\infty]=0\). For \(p=3,7\) there is a witness on every disc.
At the remaining primes the upper and lower counts coincide.
In particular, at \(p=41\) successive precision levels reduce the
upper count from \(40\) to \(12\) and then \(10\).
\begin{table}[htbp]
\centering
\caption{Classical abelian Chabauty on \(C\). In each group \(n\) counts
all residue discs and \(o\) counts the occupied discs.}\label{tab:genus}
\begin{tabular}{rrr@{\qquad}rrr@{\qquad}rrr}
\toprule
\(p\)&\(n\)&\(o\)&\(p\)&\(n\)&\(o\)&\(p\)&\(n\)&\(o\)\\
\midrule
3&5&5&59&61&5&131&161&5\\
5&6&4&61&62&8&137&140&10\\
7&10&10&67&69&5&139&138&8\\
11&16&6&71&79&9&149&164&6\\
13&15&5&73&72&10&151&140&12\\
17&15&11&79&75&11&157&152&6\\
19&28&6&83&72&6&163&165&7\\
23&30&10&89&81&9&167&139&9\\
29&30&6&97&103&9&173&176&6\\
31&29&9&101&106&6&179&183&5\\
37&39&5&103&103&9&181&185&5\\
41&40&10&107&96&6&191&201&11\\
43&47&7&109&109&5&193&208&10\\
47&45&9&113&124&10&197&211&5\\
53&62&6&127&145&9&199&184&10\\
\bottomrule
\end{tabular}
\end{table}

In the holomorphic basis the Cartier--Manin matrix has entries
\([x^{pi-j}]f^{(p-1)/2}\), \(1\le i,j\le3\), up to the transpose
convention \cite{AchterHowe}. Its rank is three at 43 of these primes
and two at 5 and 13. Since Cartier annihilates exact differentials, a
nonzero reduced annihilating differential at an ordinary prime cannot
have a rational primitive on the special fibre. The tangent slopes in
the test are algebraic; recovering the constants across discs calls for
further information about the lifts. Obtaining a bounded-degree relation
among those constants as \(p\) varies remains open.

\paragraph{Formalisation and reproducibility.}
The companion source is available on GitHub.
\begin{quote}
\url{https://github.com/prajaybasu/chabauty-kim-disc-bounds}
\end{quote}
It contains a single Lean package. Its logarithmic modules check the last-pivot row, adjugate denominator, bounded-height
horizontal classification, jet differentiation, and the dimension and
kernel calculations used for the lower bound. The choice of the cyclic
rows, their degree growth and the realisation by iterated integrals are
given in Section~\ref{sec:zero}. The Selmer and disc modules verify
bounded-degree elimination, primitive reduction and inverse translation.
Localisation, the Coleman comparison and the auxiliary-polynomial count
are supplied by the cited results and Section~\ref{sec:discs}.
The exceptional-prime modules verify the local tail and finite-field
implications with the stated analytic inputs. The exact tables are
checked by arithmetic programs.

\begin{samepage}
To reproduce the rank computation, run the following command.\par\nopagebreak
\begin{verbatim}
sage -python mixed_rank.py --certify 1 2 3
\end{verbatim}
\end{samepage}
To reproduce the 45-prime census and its witnesses, run this command in
the same directory.
\begin{verbatim}
sage -python genus_three_jacobian.py
\end{verbatim}
